% ---------------------------------------------------------------- abbreviations
\section*{List of Abbreviations}
\addcontentsline{toc}{section}{List of Abbreviations}

\begin{longtable}{@{}p{0.18\textwidth}p{0.76\textwidth}@{}}
\toprule
\textbf{Abbreviation} & \textbf{Meaning} \\
\midrule
\endhead
AMP      & Automatic Mixed Precision \\
ASR      & Automatic Speech Recognition \\
AVSR     & Audio-Visual Speech Recognition \\
BiLSTM   & Bidirectional Long Short-Term Memory network \\
BN       & Batch Normalisation \\
CER      & Character Error Rate \\
CNN      & Convolutional Neural Network \\
CTC      & Connectionist Temporal Classification \\
DCT      & Discrete Cosine Transform \\
FPS      & Frames Per Second \\
GPU      & Graphics Processing Unit \\
LM       & Language Model \\
LSTM     & Long Short-Term Memory \\
MFCC     & Mel-Frequency Cepstral Coefficients \\
MHA      & Multi-Head Attention \\
MUSAN    & Music, Speech and Noise corpus \\
RMS      & Root Mean Square \\
RQ       & Research Question \\
RTF      & Real-Time Factor \\
SNR      & Signal-to-Noise Ratio \\
STFT     & Short-Time Fourier Transform \\
UI       & User Interface \\
WER      & Word Error Rate \\
\bottomrule
\end{longtable}

\subsection*{Statement labels used in this document}
To separate fixed project requirements from engineering choices, statements
are tagged as follows:

\begin{description}[leftmargin=3.2cm,style=nextline]
  \item[\spec] Specified by the LipScribe project proposal. These items define
        the intended system and are not changed without team agreement.
  \item[\rec] An implementation decision recommended in this document. It is
        a practical starting point and may be revised on evidence.
  \item[\tbd] An open question that will be settled experimentally or by the
        team during development.
\end{description}
